\reading{LTR-14}{Liquidity Transfer Pricing: A Guide to Better Practice}{STRIKE FIRST}{2}

\begin{lrkeybox}[Learning objectives — official 2026 spine wording]
\begin{enumerate}[label=\textbf{\alph*.}]
\item Discuss the process of liquidity transfer pricing (LTP) and identify best practices for the governance and implementation of an LTP process.
\item Discuss challenges that may arise for banks during the implementation of LTP.
\item Compare the various approaches to liquidity transfer pricing (zero cost, average cost, and matched-maturity marginal cost).
\item Describe the contingent liquidity risk pricing process and calculate the cost of contingent liquidity risk.
\end{enumerate}
{\footnotesize\color{FRMGrey}Source: Joel Grant, ``Liquidity Transfer Pricing: A Guide to Better Practice,'' Occasional Paper, Financial Stability Board / BIS, 2011.}
\end{lrkeybox}

% =====================================================================
\lo{LTR-14 a}{Discuss the process of liquidity transfer pricing (LTP) and identify best practices for the governance and implementation of an LTP process.}

\begin{lrdefbox}[Liquidity transfer pricing]
A process that attributes the \textbf{costs, benefits and risks of liquidity} to the respective business units within a bank. Its purpose is to transfer liquidity costs and benefits from business units to a \textbf{centrally managed pool}.
\end{lrdefbox}

\begin{lrkeybox}[The two directions of the transfer]
To achieve this, LTP:
\begin{itemize}
\item \textbf{charges users of funds} --- assets and loans --- for the \textbf{cost} of liquidity, and
\item \textbf{credits providers of funds} --- liabilities and deposits --- for the \textbf{benefit} of liquidity.
\end{itemize}
\end{lrkeybox}

\begin{lrtrapbox}[Charge the asset side, credit the liability side]
This is a pure polarity pair and it is the first thing to lock. \textbf{Assets consume} liquidity and are \emph{charged}; \textbf{liabilities supply} it and are \emph{credited}. A distractor reversing the two reads perfectly well to anyone who has not said the direction out loud.
\end{lrtrapbox}

% =====================================================================
\lo{LTR-14 b}{Discuss challenges that may arise for banks during the implementation of LTP.}

Banks historically \textbf{mispriced liquidity}, leading to an over-reliance on short-term funding for long-term assets. Five challenges follow:

\begin{enumerate}
\item \textbf{Oversight.} Better oversight is needed, to ensure banks charge enough for long-term loans and credit deposits fairly.
\item \textbf{Centralization.} A \textbf{centralized treasury} can help prevent mispricing and keep the whole bank on the same basis.
\item \textbf{Trading activities.} Clear rules are needed for trading activities, to account for liquidity risks.
\item \textbf{Systems.} The \term{Liquidity Management Information System (LMIS)} needs upgrading, to track liquidity costs and benefits accurately.
\item \textbf{Incentives.} Bonuses and rewards should not focus on profit alone, but also on how risky an activity is \emph{in liquidity terms}.
\end{enumerate}

\begin{lrtrapbox}[The incentive point is the one that gets dropped]
Four of the five challenges are about measurement and process. The fifth --- \textbf{compensation must reflect liquidity risk, not just profit} --- is about behaviour, and it is the one an incomplete option leaves out. It is also the mechanism that connects LTP to the conduct failures the crisis exposed.
\end{lrtrapbox}

% =====================================================================
\lo{LTR-14 c}{Compare the various approaches to liquidity transfer pricing (zero cost, average cost, and matched-maturity marginal cost).}

\begin{center}
\scriptsize
\begin{longtable}{@{}p{3.4cm} p{5.7cm} p{5.2cm}@{}}
\toprule
\rowcolor{FRMSoft}\textbf{Approach} & \textbf{How it works} & \textbf{Problem} \\
\midrule
\endfirsthead
\toprule
\rowcolor{FRMSoft}\textbf{Approach} & \textbf{How it works} & \textbf{Problem} \\
\midrule
\endhead
\textbf{1. Zero cost} &
Banks do not account for the costs, benefits and risks associated with liquidity. Liquidity is treated \textbf{as if it were free}. &
Balance sheets become heavily skewed towards \textbf{illiquid assets without sufficient long-term liabilities} to balance the funding need. During the crisis the flaws became evident as liquidity dried up and the cost of funding soared. \\
\addlinespace
\textbf{2. Pooled average cost of funds} &
Average the interest expense across \textbf{all funding sources and maturities}, then apply that single average cost to both assets and liabilities. &
Simple, but does not differentiate liquidity cost by maturity. It \textbf{undercharges long-term assets} --- which carry higher liquidity risk --- and \textbf{overcharges short-term assets}. It likewise fails to differentiate the credit given to liabilities by duration. \\
\addlinespace
\textbf{3. Separate average costs of funds} &
Uses \textbf{different average costs for assets and for liabilities}, potentially varying the spreads between them. Aims to better align costs and benefits with their respective maturities. &
Allows some differentiation, but \textbf{still uses averages} that may not reflect current market conditions or the specific risks of different maturities. \\
\addlinespace
\textbf{4. Matched-maturity marginal cost of funds} &
The most sophisticated and accurate method. Aligns the cost of funds with the \textbf{actual market cost of funding} and the \textbf{specific maturities} of assets and liabilities, using the bank's actual cost of funding to calculate the correct liquidity spread. &
--- \textbf{Advantages:} a more precise measure, reflecting current market conditions and maturity-specific risks. Particularly useful in \textbf{volatile market conditions}, when liquidity costs can fluctuate significantly. \\
\bottomrule
\end{longtable}
\end{center}
\figcap{The spine names three approaches; the source notes give four, splitting the average-cost method into pooled and separate variants. The sequence is a single axis of increasing granularity: no differentiation, one average for everything, separate averages by side, then marginal cost matched to each maturity.}

\begin{lrtrapbox}[The undercharge and overcharge run in specific directions]
Pooled average cost \textbf{undercharges long-term assets} and \textbf{overcharges short-term assets}. That is not arbitrary: a single average sits below the true cost of long funding and above the true cost of short funding. Reversing the pair is the polarity build, and it is more tempting than it looks because ``averaging overcharges the risky thing'' feels intuitive and is wrong.
\end{lrtrapbox}

% =====================================================================
\lo{LTR-14 d}{Describe the contingent liquidity risk pricing process and calculate the cost of contingent liquidity risk.}

\begin{lrkeybox}[What the source notes carry]
Calculating the charge for using --- or the credit for providing --- funding liquidity for \textbf{contingent commitments} such as lines of credit, collateral postings for derivatives and other financial contracts, is quite difficult. In these cases the best approach is to \textbf{impose a scenario model}. Banks carry a liquidity cushion, or \textbf{buffer}, to help them survive periods of unexpected funding outflows.
\end{lrkeybox}

\gapbadge
\gapnote{The source notes describe the process and stop. The formula and worked example below --- the ``calculate'' half of the objective --- are written from the GARP curriculum (Grant, Sections 15.4.6 and 15.4.7).}

\begin{lrgapbox}
The objective says \emph{describe and calculate}. The described half is present; the computational half, which turns a contingent commitment into a basis-point charge, is not.
\end{lrgapbox}

\begin{lrfmlbox}[Funds transfer price, and the cost of contingent liquidity risk]
\[ \text{FTP} \;=\; \text{base rate} \;+\; \text{term liquidity premium} \;+\; \text{liquidity premium} \]
\vk{\textbf{base rate} = the rate from the swap curve corresponding to the asset's contractual or behavioural maturity (or repricing term, whichever is less); \textbf{term liquidity premium} = the spread between the swap curve and the bank's marginal cost of funds curve, based on that same maturity; \textbf{liquidity premium} = the cost of carrying the liquidity cushion, averaged over the bank's total assets.}
\vspace{4pt}
\begin{align*}
\text{Cost of contingent liquidity risk} \;=\;& \frac{\text{Limit} - \text{Drawn amount}}{\text{Limit}} \;\times\; \text{Likelihood of drawdown} \\[2pt]
&\times\; \text{Cost of funding the liquidity cushion}
\end{align*}
\vk{The first term is the \emph{undrawn proportion} of the facility. The \term{likelihood of drawdown} (or drawdown factor) is assessed using \textbf{behavioural modelling}, and depends on factors such as the customer's drawdown history, the customer's credit rating, and whatever else the bank deems relevant.}
\end{lrfmlbox}

\begin{lrexbox}[Pricing an undrawn line of credit]
A line of credit has a limit of \$10 million with \$4 million already drawn. Behavioural modelling puts the likelihood the customer draws on the remaining credit at \textbf{60 per cent}. The cost of term funding the assets held in the liquidity cushion is \textbf{18 basis points}, taken from the three-year term liquidity premium. What rate should be charged for contingent liquidity risk?
\tcblower
\textbf{Step 1 --- the undrawn proportion of the facility.}
\[ \frac{\$10\text{m} - \$4\text{m}}{\$10\text{m}} = 0.60 \]
Note that the \emph{drawn} \$4 million is not charged here --- it is already an on-balance-sheet asset carrying its own FTP charge. Only the \textbf{contingent} part is being priced.

\textbf{Step 2 --- apply the drawdown factor.}
\[ 0.60 \times 0.60 = 0.36 \]
This is the expected proportion of the \emph{limit} that will become a funded exposure.

\textbf{Step 3 --- apply the cost of funding the cushion.}
\[ 0.36 \times 0.0018 = 0.000648 = \mathbf{6.48\text{ basis points}} \]

\textbf{What this means.} The bank must hold liquid assets against the possibility of the draw, and holding them has a carrying cost. 6.48 bp is that carrying cost pushed back to the business unit, product or transaction that created the need. Without this charge, contingent commitments look free to the desk writing them --- which is precisely why they were neglected before the crisis.
\end{lrexbox}

\begin{lrtrapbox}[Three ways the 6.48 bp is corrupted]
\begin{itemize}
\item \textbf{Charging the whole limit.} $1.00 \times 0.60 \times 0.0018 = 10.8$ bp. The drawn portion is excluded; only the undrawn proportion is contingent.
\item \textbf{Charging the drawn portion.} $0.40 \times 0.60 \times 0.0018 = 4.32$ bp --- the exact complement, and it reads as plausible.
\item \textbf{Dropping the drawdown factor.} $0.60 \times 0.0018 = 10.8$ bp, the same number as the first error by coincidence of this data. Two different mistakes converging on one offered option is exactly the \S5B.3 shape.
\end{itemize}
\end{lrtrapbox}

\begin{lrkeybox}[Why contingent liquidity risk went unpriced, and what better practice requires]
The uncertainty surrounding future cash-flow demands from contingent commitments makes them particularly difficult to assess and price --- one reason they were neglected before the GFC. Products that received little attention and then required significant funding included \textbf{credit card loans and investments, trading positions and derivatives, revolving lines of credit, and liquidity lines}.

The first step towards better management is \emph{not} to ask how much should be charged, but for banks to understand that \textbf{all contingent commitments need to be charged}. Once that is settled, the pricing methods can be refined. A further problem with not having a \textbf{granular} charge is that it encourages businesses to deal in products that are not being charged for the contingent liquidity risk they actually present --- the behaviour-distorting effect of under-pricing.
\end{lrkeybox}

\begin{lrtrapbox}[Consolidated trap summary — LTR-14]
\begin{itemize}
\item \textbf{Polarity.} LTP \textbf{charges} users of funds (assets, loans) and \textbf{credits} providers of funds (liabilities, deposits).
\item \textbf{Polarity.} Pooled average cost \textbf{undercharges long-term assets} and \textbf{overcharges short-term assets}.
\item \textbf{Sibling.} Zero cost $\rightarrow$ pooled average $\rightarrow$ separate averages $\rightarrow$ matched-maturity marginal. One axis, increasing granularity; matched-maturity marginal is the most accurate and the most useful in volatile markets.
\item \textbf{Calculation.} Cost of contingent liquidity risk $=$ \textbf{undrawn} proportion $\times$ drawdown likelihood $\times$ cost of funding the cushion. The drawn portion is excluded.
\item \textbf{Definition.} The drawdown likelihood comes from \textbf{behavioural modelling} --- drawdown history, credit rating --- not from a regulatory factor.
\item \textbf{Scope.} FTP $=$ base rate $+$ term liquidity premium $+$ liquidity premium, where the last term is the \textbf{cost of carrying the cushion averaged over total assets}.
\item \textbf{Scope.} The first step in better practice is recognising that \textbf{all} contingent commitments must be charged --- not determining the amount.
\item \textbf{Scope.} Incentive compensation must reflect liquidity risk, not profit alone.
\end{itemize}
\end{lrtrapbox}
